% One methods section for a journal manuscript.
% Paste from \section through the end of the section, or compile this file alone.
\documentclass[11pt]{article}
\usepackage[margin=1in]{geometry}
\usepackage{amsmath,amssymb}
\usepackage{booktabs}
\usepackage{microtype}

\begin{document}

\section{Software}
\label{sec:software}

The calculations reported here are performed with a physics-informed neural
network for the one-dimensional shallow-water equations. A single trainer
builds a two-field network for water depth and velocity, minimizes a
strong-form residual, and, when a search is requested, selects the
architecture with Optuna. The differential equation, the initial and
boundary data, and any exact reference solution are supplied by a test-case
object. The trainer does not contain a separate discretization for each
benchmark. All trainable tensors and automatic derivatives are stored in
64-bit precision. The implementation uses TensorFlow and Keras for the
network, the optimizer, and the spatial and temporal derivatives; NumPy for
collocation coordinates and archived fields; Optuna for the hyperparameter
study; and SciPy for optional L-BFGS-B refinement.

\subsection{Network families}
\label{sec:families}

The network returns the pair of fields supplied to the test case, one
tower for each field. Hidden layers are
dense maps with a hyperbolic-tangent activation and Glorot normal
initialization. Output units are linear. Three families share this
interface.

In the classical family, each tower is a multilayer perceptron with \(d\)
hidden layers of width \(w\). For one tower with a two-dimensional input
and \(d\ge 1\), the number of trainable parameters is
\begin{equation}
\label{eq:classical-count}
3w + (d-1)(w^{2}+w) + (w+1).
\end{equation}
The model count is twice~\eqref{eq:classical-count}. No circuit is
constructed.

In the hybrid family, those perceptrons are retained and a quantum filter
is added on each tower after the last hidden layer,
\begin{equation}
\label{eq:hybrid}
z \leftarrow z + \alpha\,
\operatorname{decode}\!\big(\operatorname{circuit}(\operatorname{encode}(z))\big).
\end{equation}
The encoder is a hyperbolic-tangent map from the trunk width onto one
feature per qubit. The decoder maps the measured expectations back to the
trunk width. The gate \(\alpha\) is a trainable scalar, initialized at
zero, so the hybrid model begins as its classical trunk. Each tower has
its own encoder, circuit, decoder, and gate. A hybrid specification with
zero qubits is stored and trained as a classical model.

In the full-quantum family there is no hidden perceptron and no residual
gate. Each field is
\begin{equation}
\label{eq:fullq}
\operatorname{readout}\!\big(\operatorname{circuit}(\operatorname{encode}(x,t))\big).
\end{equation}
The encoder is one hyperbolic-tangent map from the two coordinates onto
the qubit features, and the readout is a linear map from the expectations
to a scalar. Width and depth are not used. The encoder and readout are
the parameters counted as classical for this family. With \(n_{q}\) qubits
on each of the two towers, that count is \(8n_{q}+2\).

Both quantum families use the same circuit. A layer applies a three-angle
rotation on every qubit and then a ring of controlled-NOT gates, including
the gate from the last qubit back to the first. The feature vector enters
through an angle embedding of \(Y\) rotations, and the circuit returns the
Pauli-\(Z\) expectation of every qubit. The number of rotation angles on
one tower is \(3n_{q}n_{L}\), so the two towers contribute
\(6n_{q}n_{L}\) quantum parameters. The statevector is evaluated with
PennyLane's \texttt{default.qubit} device, batched and differentiated
through a JAX just-in-time compilation. TensorFlow retains the rotation
angles and the residual tape. A custom gradient inserts the JAX
vector-Jacobian product into that tape. The compiled circuit is built
again whenever the qubit count or the layer count changes.

A hybrid run may disable the circuit for an initial fraction of the Adam
iterations. During that warmup the forward evaluation is the classical
trunk, while the quantum variables remain in the optimizer and receive
zero gradients. The full-quantum family has no trunk on which to perform
that warmup, so the fraction is ignored there.

\subsection{Residual loss}
\label{sec:loss}

The loss is the sum of the strong-form residual of the governing equations
and the initial and boundary data. At the collocation stations the test
case returns the two residual components \(r_{m}\) and \(r_{p}\). With
logarithmic weights \(\lambda_{m}\) and \(\lambda_{p}\), the normalized
coefficients of those components are
\begin{equation}
(w_{m},w_{p})
=
\frac{\big(e^{\lambda_{m}},e^{\lambda_{p}}\big)}
{\operatorname{mean}\big(e^{\lambda_{m}},e^{\lambda_{p}}\big)},
\end{equation}
so equal logarithms produce unit weights. Each residual contribution is a
mean square. Derivatives of the network outputs are obtained by
reverse-mode automatic differentiation.

The term \(L_{\mathrm{soft}}\) is the initial and boundary data. It is the
weighted sum
\begin{equation}
\label{eq:lsoft}
L_{\mathrm{soft}} = L_{\mathrm{ic}} + L_{\mathrm{bc}},
\end{equation}
where \(L_{\mathrm{ic}}\) is the mean-square mismatch between the network
and the prescribed initial values, and \(L_{\mathrm{bc}}\) is the
mean-square mismatch with the prescribed boundary values. The coefficients
on these data belong to the test case. They are independent of \(w_{m}\)
and \(w_{p}\). The discrete loss is
\begin{equation}
\label{eq:loss}
L = w_{m}\,\operatorname{mean}(r_{m}^{2})
  + w_{p}\,\operatorname{mean}(r_{p}^{2})
  + L_{\mathrm{soft}}.
\end{equation}
In the soft-constraint mode the network output is free, and
\eqref{eq:lsoft} enforces the initial and boundary data. In the
hard-constraint mode those data are built into the network output, so
they hold identically and \(L_{\mathrm{soft}}=0\). Adam and L-BFGS-B
minimize~\eqref{eq:loss}.

Collocation points form one fixed tensor product on the test-case
rectangle and are held constant for the whole optimization. Each lattice
is specified by its number of space stations \(n_{x}\) and its number of
time stations \(n_{t}\). Both axes include their endpoints, and the
number of collocation points is the product \(n_{x}n_{t}\). The trial,
the objective evaluation, the final Adam run, and L-BFGS-B each have
their own pair \((n_{x},n_{t})\). Adam uses a constant learning rate.
If L-BFGS-B is enabled, it minimizes the same loss on its lattice, with
the TensorFlow gradient supplied to SciPy as the Jacobian.

An Optuna trial trains for a prescribed number of Adam epochs on its
collocation lattice. The value recorded for that trial is~\eqref{eq:loss}
evaluated once more, without a further parameter update, on the objective
lattice. That value includes both the governing residual and
\(L_{\mathrm{soft}}\). Comparison with an exact depth or velocity is
performed only after the selected model has been trained again, on a
separate plotting mesh.

The study may minimize \(L\) itself, the product of \(L\) with a factor
linear in the trainable-parameter count, or a parameter count subject to
a prescribed ceiling on \(L\). In every mode the parameter count, the
split between circuit angles and the remaining weights, and the separate
residual and penalty contributions are stored with the trial.

\subsection{Hyperparameter search}

Optuna minimizes one scalar. The categorical choice of family is made
first. Width and depth are proposed only for the classical and hybrid
families. Qubit count and circuit depth are proposed only for the hybrid
and full-quantum families. Learning rate is drawn log-uniformly. The
logarithmic loss weights, when they are included in the search, are drawn
uniformly. Family-specific bounds may be placed on width and depth. Those
bounds may be set equal, so a hybrid trial and a classical trial search
the same trunk and differ by the residual circuit~\eqref{eq:hybrid}.

The sampler is the tree-structured Parzen estimator, with grouped
conditional parameters. A prescribed number of initial trials is drawn
uniformly. Setting that number equal to the trial budget keeps every draw
uniform, and each trial then selects a family independently. The budget
is a cap on completed trials in one SQLite study. Independent processes
share that study and stop at the same cap. The seed of process \(k\) is
\(42+k\).

After the search, the winning family, widths, learning rate, circuit
size, and loss weights are used to initialize a new network. That network
is trained from the start for the final Adam budget. The weights of the
winning trial are not reused. The final approximation is the field
produced by this second training.

\subsection{Test-case interface}
\label{sec:interface}

A test case is a builder that returns a domain rectangle, an output map,
a residual map, and, where they exist, an exact solution and a soft
penalty. The output map receives the raw coordinates and the constraint
mode, forms the two-column network input, and returns the two fields.
In the hard mode that map composes the network output so the prescribed
initial and boundary values hold identically. In the soft mode the
network output is unconstrained, and the penalty \(L_{\mathrm{soft}}\)
measures the mismatch with the prescribed data under coefficients
belonging to the test case. The residual map differentiates the returned
fields and returns \((r_{m},r_{p})\). The trainer then
evaluates~\eqref{eq:loss}. An exact solution, when the builder provides
one, is used for the final comparison on the plotting mesh. It enters
the loss only if the test case places it inside \(L_{\mathrm{soft}}\).

Adding a test case consists of defining its physical constants,
implementing this builder, and registering the builder under a new name.
The search database for that test case is kept separate, because a stored
study freezes the distributions from which its trials were drawn and
because its objective belongs to a different residual.

\end{document}
